% Generated draft: requires mathematical and source review.
\documentclass[11pt]{article}
\usepackage[margin=25mm]{geometry}
\usepackage{fontspec}
\setmainfont{Latin Modern Roman}
\usepackage{amsmath,amssymb,mathtools,unicode-math}
\setmathfont{Latin Modern Math}
\usepackage{enumitem,needspace,xurl,hyperref}
\setlength{\parindent}{0pt}
\setlength{\parskip}{5pt}
\setlength{\emergencystretch}{3em}
\newcommand{\N}{\mathbb{N}}
\newcommand{\Z}{\mathbb{Z}}
\newcommand{\Q}{\mathbb{Q}}
\newcommand{\R}{\mathbb{R}}
\newcommand{\C}{\mathbb{C}}
\newcommand{\F}{\mathbb{F}}
\newcommand{\A}{\mathbb{A}}
\newcommand{\PP}{\mathbb{P}}
\newcommand{\E}{\mathbb{E}}
\newcommand{\eps}{\varepsilon}
\newcommand{\dd}{\,\mathrm{d}}
\begin{document}
\section*{UnsolvedMath ID: 30001552}
\textbf{Source problem code:} catalogue entry
\subsection{Definitions and mathematical statement}
Let \(A\) be an alphabet, and let \(A^*\) denote the set of finite words over \(A\), including the empty word \(\varepsilon\), with concatenation as its operation. For \(w\in A^*\), write \(|w|\) for its length. An antimorphic involution is a map \(\theta:A^*\to A^*\) satisfying
\[
\theta(uv)=\theta(v)\theta(u)\quad(u,v\in A^*),
\qquad
\theta(\theta(u))=u\quad(u\in A^*).
\]
For integer lengths \(p,q\) for which it is defined, \(\gcd(p,q)\) denotes their greatest common divisor.

The supplied question is: for every such \(\theta\), every \(w\in A^*\), and every pair \(p,q\) of alternating \(\theta\)-periods of \(w\), does the implication
\[
|w|\geq p+q-\gcd(p,q)
\quad\Longrightarrow\quad
\gcd(p,q)\text{ is a }\theta\text{-period of }w
\]
hold?

The supplied record defines neither an alternating \(\theta\)-period nor a \(\theta\)-period. These two terms have therefore been retained as distinct source terms without assigning them an interpretation. The record also does not specify the admissible ranges of period lengths, including positivity, whether lengths exceeding \(|w|\) are allowed, or conventions for the empty word. It imposes no explicit finiteness condition on \(A\). Consequently, this is a transcription of the intended universal implication with unresolved predicates and boundary conventions, rather than a complete standalone formalization.
\subsection{Short English statement}
The question asks whether two alternating \(\theta\)-periods \(p,q\) force their greatest common divisor to be a \(\theta\)-period whenever the word has length at least \(p+q-\gcd(p,q)\). The supplied definitions of the two period notions are missing, so their intended meanings and boundary conventions require review.
\subsection{Sources}
\begin{itemize}
\item Supplied local UnsolvedMath registry record, permanent ID \(30001552\), title: Fine--Wilf Bounds for Antimorphic Periods; collection: Oberwolfach Reports Open Problems; exact source problem code: \texttt{catalogue entry}. No individual proposer is identified. The supplied proposed year is 2010.
\item Supplied source citation: Mini-Workshop: Combinatorics on Words (2011), Oberwolfach Reports, pp.\ 2195--2244, DOI \texttt{10.4171/owr/2010/37}. The source URL and the literature-source link labelled ``original workshop report'' are both \url{https://doi.org/10.4171/owr/2010/37}. The supplied background describes Oberwolfach Reports as documenting mathematical workshops and their open-problem discussions. The citation's 2011 date and the record's 2010 year attribution require reconciliation.
\item Literature assessment dated exactly \texttt{2026-08-21}: the supplied curation labels the problem open and reports that the sharper alternating-antimorphic bound involving the greatest common divisor remains a conjecture, with only weaker bounds proved. This is supplied dated curation, not an independently verified assessment. The background explicitly requests independent verification of time-sensitive status and cited claims.
\item The supplied statement-status label is \texttt{corrected\_verified}; its accompanying verification note says only that TeX notation was enclosed in math delimiters without changing the mathematical content. This does not establish mathematical correctness or resolve the missing definitions.
\end{itemize}
\end{document}
